\scanpage{150}
\[
\tag{57}
\begin{pmatrix}
 A+2\unsure{\nu}C&B+2\unsure{\nu}D\\
 C&D
\end{pmatrix}
=\varepsilon\begin{pmatrix}R&S\\T\,(=S)&U\end{pmatrix}.
\]

But this condition implies first that the matrix on the left-hand side is symmetric (since $S=T$) i.e.
\[
2\nu D=C-B
\]
a condition which is also equivalent, by (54), to
\[
\tag{58}
\inserted{2\nu D=1-D\quad\text{i.e.}}\qquad
\boxed{\mu D=1}\quad
\deleted{\mu D=\unsure{1}\;\illegible{1}\qquad\unsure{\mu}=\unsure{1}}\quad
\text{i.e.}\quad D=1/\mu,\quad\mu=1/D
\]
\inserted{where $\mu=2\nu+1$,}
which already implies \uline{$\mu$ and $D$ invertible}.

\deleted{On the other hand, we must have}
Symmetrically, if we write (57) in the form
\[
\varepsilon\begin{pmatrix}A&B\\C&D\end{pmatrix}
=\unsure{\varepsilon}\begin{pmatrix}R-2\nu T&S-2\nu U\\T&U\end{pmatrix},
\]
the condition $C+D-B=1$ becomes, for the entries of the matrix $\varphi$
\[
T+U-(S-2\nu U)=\varepsilon\ ;
\]
i.e.
\[
(2\nu+1)U=\varepsilon\quad\text{i.e.}
\]
\[
\tag{59}
\mu U=\varepsilon\quad\text{i.e.}\quad U=\varepsilon/\mu,
\quad\inserted{\mu=\varepsilon/U}
\]
which is also equivalent, taking (58) into account, to the relation
\[
\tag{60}
D=\varepsilon U\quad\text{or}\quad U=\varepsilon D
\]
which is none other than the equality of the entries in the bottom rows of the matrices (57).
\margin{149}

\scanpage{151}
\mspage{365}
We can write out the relations (57), as \deleted{\illegible{1}} four relations on the $9$ variables $a,b,c,d,r,s,t,u,\mu$ in addition to the relations (48), (49);
\[
\tag{61}
\left\{\begin{gathered}
\boxed{\begin{aligned}
D&=\varepsilon U\\
\mu D&=1\\
C&=\varepsilon T
\end{aligned}}
\quad\begin{aligned}
\deleted{\mu\unsure{U}=\unsure{\varepsilon}}{i.e.}\quad\deleted{\unsure{D=1/\mu}}\quad\mu U=\varepsilon\\
&
\end{aligned}\\
A=\varepsilon\bigl(R-\deleted{\unsure{2\nu}}\underbrace{(\mu-1)}_{2\nu}\underbrace{T}_{C}\bigr)
\end{gathered}\right.
\]
\margin{\uline{N.B.} this last one is a consequence of the others, taking $C+D-B=1$ and $AD-BC=1$ into account}
i.e. -- \deleted{or} \inserted{again}
\[
\tag{62}
\left\{\begin{gathered}
\boxed{\begin{aligned}
\mu(c^2+d^2+cd)&=\delta\\
\unsure{c}^2+d^2+cd&=\varepsilon\delta(t^2+u^2)\\
\uline{ac}+\uline{bd}+\uline{ad}-(c^2+d^2+cd)&=\delta\varepsilon(rt+su)
\end{aligned}}\\
a^2+b^2+ab-(ac+bd+bc)=\delta\varepsilon\bigl(r^2+s^2-(\mu-1)(rt+su)\bigr)
\end{gathered}\right.
\]
\margin{consequence of the preceding ones}

When the signs $\varepsilon,\delta$ are fixed, we find \inserted{with (48), (49)} a system of $6$ equations in $9$ variables $a,b,c,d,r,s,t,u,\mu$ -- we should be able to check that the scheme (over $\operatorname{Spec}\mathbb Z$) of solutions of these equations is of relative dimension $3$, and of relative dimension $1$ over the ``factors'' $(a,b,c,d)$ resp. $(r,s,t,u)$. We thus have an adelic point of this absolute scheme -- a point with coefficients in $\widehat{\mathbb Z}$\ldots
\margin{no, it is of total relative dimension $4$}

\uline{Theorem} The multiplier $\mu$ is determined as a function of $\unsure{\xi|\zeta}\inserted{=\begin{pmatrix}A&B\\C&D\end{pmatrix}}$ (considered as an element of $\operatorname{Sl}(2,\widehat{\mathbb Z})$), by the relation
\[
\mu D=1\qquad(58)
\]
\unsure{\textup{(+)}}
and it is determined \uline{up to sign} by $\varphi=\begin{pmatrix}R&S\\T&U\end{pmatrix}$ by the relation
\[
\mu U=\varepsilon\qquad(59)\qquad\text{where }\varepsilon\in\{\pm1\},
\]
\margin{150}

\scanpage{152}
and where $\varepsilon$ is the sign which occurs in the formulas
\[
\left\{\begin{aligned}
\unsure{\xi|\zeta}&=\varepsilon l_0^\nu\varphi\qquad(56)\\
g\rho(g)\rho^2(g)&=\varepsilon\qquad(45).\\
\inserted{g}&\inserted{=\varepsilon\alpha\rho(\alpha)^{-1}}
\end{aligned}\right.
\]
We have $\varepsilon=1$ iff $\mu\equiv1\ (4)$, $\varepsilon=-1$ iff $\mu\equiv-1\ (4)$.
\inserted{i.e. $\nu\equiv0\ (2)$}
\inserted{i.e. $\nu\equiv1\ (2)$}

\margin{in other words, we have $\mu\equiv\varepsilon\ (4)$ since $U=t^2+u^2\equiv1\ (4)$ hence $1/U\equiv(1)\ (4)$ hence $\mu=\varepsilon/U\equiv\varepsilon\ (4)$. This would be the time to include \unsure{also the|here the} relation $\mu\equiv\delta\ (3)$ cf. with the appropriate proof\ldots}

To prove this, we recall that
\[
U=t^2+u^2
\]
note that the only values modulo $4$ taken by $t^2$ and $u^2$ are $0$ and $1$, hence those taken by $t^2+u^2$ are $0,1$ and $2$, and the only unit among these is $1$. Thus \deleted{\unsure{$\mu$}} $1/\mu$ can be put in the form $\varepsilon(\unsure{t^2}+u^2)$ modulo $4$ iff $1/\mu\equiv\varepsilon\ (4)$ or again $\mu\equiv\varepsilon\ (4)$ q.e.d.

\deleted{Thus} We had seen in a quite analogous way, thanks to the formula $\mu D=1\quad(58)$ and
\[
D=\delta(c^2+d^2+cd)
\]
that $\mu\equiv\delta\ (3)$ i.e. $\delta=+1$ (resp. $-1$) iff $\mu\equiv1\ (3)$ (resp. $\mu\equiv-1\ \mathrm{mod}\ 3$). We \deleted{\unsure{could thus}} shall denote, for \unsure{all|the} prime numbers $p$, by $\unsure{\varepsilon_p|\beta_p}$ the value at $\mu$ of the ``quadratic character'' relative to $p$, equal to $1$ iff $\mu_p\in(\mathbb Z_p^*)^2$. We shall therefore have
\[
\tag{63}
\left\{\begin{aligned}
\delta&=ad-bc=\unsure{\varepsilon_3|\beta_3}(\mu)=\varepsilon_3\\
\varepsilon&=D/U=\unsure{\varepsilon_2|\beta_2}(\mu)=\varepsilon_2,
\end{aligned}\right.
\]
and knowing $\varepsilon_3$ (resp. $\varepsilon_2$, resp. both) is equivalent to knowing $\mu$ modulo $\unsure{3|2}$ (resp. modulo $4$, resp. modulo $12$).
\margin{$(\mathbb Z/12\mathbb Z)^*\simeq(\mathbb Z/4\mathbb Z)^*\times(\mathbb Z/3\mathbb Z)^*\simeq\{\unsure{\pm1}\}\times\{\pm1\}$}
\margin{151}

\scanpage{153}
\mspage{366}
We are also going to compute \deleted{\illegible{1}} (cf. formulas (45))
\[
\alpha=(l_0^\nu f)^{-1}\alpha_0=f^{-1}l_0^{-\nu}\alpha_0
\]
and
\[
g=\varepsilon\alpha\rho(\alpha)^{-1}
=\underbrace{\sigma_\infty(f)^{-1}}_{{}^t f}\,l_1^\nu\rho(f)
=(l_0^\nu f)^{-1}\unsure{\xi|\zeta}\,\rho(l_0^\nu f).
\]
We find
\[
\tag{64}
\begin{aligned}
\alpha&=\underbrace{\begin{pmatrix}u&-s\\-t&r\end{pmatrix}}_{f^{-1}}
\underbrace{\begin{pmatrix}1&2\nu\\0&1\end{pmatrix}
\begin{pmatrix}a&b\\c&d\end{pmatrix}}_{\begin{pmatrix}a+2\nu c&b+2\nu d\\c&d\end{pmatrix}}\\
&=\begin{pmatrix}u(a+2\nu c)-cs&u(b+2\nu d)-ds\\-t(a+2\nu c)+cr&-t(b+2\nu d)+dr\end{pmatrix},
\end{aligned}
\]
\margin{\uline{N.B.} Unlike $\alpha_0$ and $f$, and $\unsure{\xi|\zeta},\varphi$ which are deduced from them, $\alpha$ \deleted{\illegible{1}} shares with $g$ the advantage and the interest of being determined without ambiguity (up to sign) by an automorphism $u$ of $\widehat\Pi_{0,3}$, whereas $f$ etc. are \illegible{1} \unsure{choices of ambiguities} $f\mapsto l_0^n f$ etc.\ldots}
\[
\begin{aligned}
g&=\underbrace{\begin{pmatrix}r&t\\s&u\end{pmatrix}}_{\sigma_\infty(f^{-1})={}^t f}
\underbrace{\begin{pmatrix}1&0\\2\nu&1\end{pmatrix}}_{\unsure{l_1^\nu|l_0^\nu}}
\underbrace{\begin{pmatrix}t+u&-t\\\unsure{t+u-r-s}&\unsure{r-t}\end{pmatrix}}_{\rho(f)\ (\text{cf. \S43 formulas }(77))}\\
&=\begin{pmatrix}r&t\\s&u\end{pmatrix}
\begin{pmatrix}t+u&-t\\\unsure{\mu(t+u)-r-s}&-\mu t+r\end{pmatrix}\\
&=\begin{pmatrix}1+\mu t(t+u)&-\mu t^2\\-1+\mu u(t+u)&1-\mu ut\end{pmatrix},
\end{aligned}
\]
i.e.
\[
\tag{65}
g=\begin{pmatrix}1+\mu t(t+u)&-\mu t^2\\-1+\mu u(t+u)&1-\mu ut\end{pmatrix}.
\]

It appears that $g$ is expressed simply in terms of $f$ i.e. of $r,s,t,u$ directly, and not in terms of $\alpha_0$ or $\unsure{\xi|\zeta}$ i.e. via the $a,b,c,d$ or $A,B,C,D$. On the other hand, to express $\alpha$ (64), we must bring in both \inserted{the} $r,s,t,u$ and the $a,b,c,d$.

It was natural, in connection with the canonical homomorphism
\[
\mathcal N_{1,1}\longrightarrow\operatorname{Gl}(2,\widehat{\mathbb Z})
\]
(cf. \S39, (11)) to look for a matrix \deleted{\illegible{1}}, which by abuse of language I am going to denote by the same letter, $u$, as the automorphism under consideration of $\widehat\Pi_{0,3}$ with which we had started (just as I denote by the same letter the elements $g,f$ etc. of $\widehat\Pi_{0,3}$ or of $\unsure{\widehat{\mathfrak S}_{0,3}|\widehat{\mathcal S}_{0,3}}$, and their images in $\operatorname{Gl}(2,\widehat{\mathbb Z})'$, or even corresponding elements in $\operatorname{Gl}(2,\widehat{\mathbb Z})$ itself),
\margin{152}

\scanpage{154}
\[
u\in\operatorname{Gl}(2,\widehat{\mathbb Z})
\]
\deleted{\unsure{satisf}} with determinant equal to $\mu$, and inducing on $\operatorname{Sl}(2,\widehat{\mathbb Z})'$, by an inner automorphism, the same automorphism as \inserted{the} $u$ \inserted{at the outset} (by passage to the quotient), which is made explicit by the formulas
\[
\tag{66}\det u=\mu
\]
\[
\tag{67}u\sigma_\infty u^{-1}=\varepsilon'\sigma_\infty\qquad\varepsilon'\in\{\pm1\}
\]
\[
\tag{68}u\rho u^{-1}=\varepsilon''g\rho\qquad\varepsilon''\in\{\pm1\}
\]
In fact, \deleted{in} (68) \deleted{we \unsure{should have}} \inserted{is written}
\[
[u,\rho]=u\rho(u)^{-1}=\varepsilon''g
\]
and applying to both sides the operation $x\mapsto x\rho(x)\rho^2(x)$, and taking into account that $g\rho(g)\rho^2(g)=\varepsilon_2$, \deleted{we find} \inserted{and $x\rho(x)\rho^2(x)=1$ for $x=u\rho(u)^{-1}$,} we find
\[
\tag{68 bis}\varepsilon''=\varepsilon_2
\]
and (68) is thus equivalent to
\[
u\rho(u)^{-1}=\varepsilon_2g=\alpha\rho(\alpha)^{-1}
\]
\deleted{\unsure{whence}} i.e. to
\[
\tag{69}u=\alpha u_0\quad\text{with }[u_0,\rho]=1
\quad\bigl(\text{i.e. }u_0\rho(u_0)^{-1}=1\bigr).
\]

We have already remarked that (67), (68) determine $u$ only modulo multiplication by a scalar $\lambda\in\widehat{\mathbb Z}^*$, and taking (66) into account, $u$ is determined only modulo $\lambda$ satisfying $\lambda^2=1$ i.e. modulo ${}_2\widehat{\mathbb Z}^*\simeq\{\pm1\}^{\mathbb P}$ (where $\mathbb P$ is the set of prime numbers), but we nevertheless hope to hit upon a canonical choice.

In my first calculations, I hadn't yet \unsure{brought out|dug up} (69), where the condition $u_0$ commutes with $\rho$ was expressed in the form
\[
\tag{70}u_0=p\,\unsure{\mathrm{id}}+q\rho=\begin{pmatrix}p&q\\-q&p+q\end{pmatrix}\quad(p,q\in\widehat{\mathbb Z}),
\]
\margin{153}

\scanpage{155}
\mspage{367}

so that
\[
\det u_0=p^2+pq+q^2
\tag{71}
\]
and that (66) is written (taking (69) into account, \unsure{where}
$\det\unsure{\alpha}=\delta=\varepsilon_3$)
\[
p^2+pq+q^2=\mu\varepsilon_3
\tag{\unsure{72}}
\]
which, in view of the relation
\[
c^2+\unsure{cd}+d^2=\varepsilon_3/\mu
\qquad(52\text{ and }58)
\]
suggested relations of the form
\[
p=\deleted{\varepsilon}\mu c,\ q=\mu d,\quad
\text{or}\quad p=\mu d,\ q=\mu c.
\]
In fact, expressing (67) as an equation in $p,q$
(added \unsure{to|:} 72), \unsure{I found}
(up to sign of \unsure{course}
which here is inessential)
\[
p=\mu d,\ q=\mu c\qquad\text{i.e.}\qquad u_0=\mu(c\rho+d)
\tag{73}
\]
but working in the case $\varepsilon'=+1$ exclusively.
(I had thought I could conclude, a little too hastily, that the
sign $\varepsilon'$ was always $+1$, before realizing
that this was hardly compatible with class
field theory, and suspecting that we should have $\varepsilon'=+1$
iff $p\equiv1\ (4)$.) By making this rather cumbersome detour
via the $(a,b,c,d)$, I nevertheless arrived, after
many twists and turns (calculation errors etc.) at the
formula where there were no longer any $a,b,c,d$ at all,
and of a disconcerting simplicity:
\[
\uuline{u}=\mu(t\sigma_\infty+u)
=\mu\begin{pmatrix}u&t\\-t&u\end{pmatrix},
\tag{74}
\]
which manifestly corresponds to the case $\unsure{\varepsilon'_2|\varepsilon'}=1$, i.e.
the one where $\uline{u}$ \uline{commutes} with $\sigma_\infty$. As

\margin{($\Longleftrightarrow$ probably, $\varepsilon_2=1$)}
\margin{I underline the letter $u$ which denotes a matrix, to
distinguish it from the $u$ which denotes a scalar.}
\margin{154}